% The observed-data sigma-algebra and the discharge of the remaining hypotheses
% of the exact MSEP theorems. Intended to be \input into content.tex.

\chapter{The observed data: Mack's hypotheses discharged}

\section{From independence across accident years to the hypotheses}

Theorem~\ref{thm:msep} conditions on a $\sigma$-algebra $\mathcal{D}$ of observed data and carries
two named hypotheses: that $\mathcal{D}$ says nothing about accident year $i$ beyond that year's own
history up to its latest observed development year. Theorem~\ref{thm:msepTotalExact} carries a
third, Definition~\ref{def:crossFree}. Mack obtains all three from assumption (M2) in the sentence
``because of the independence of the accident years'' in the proof of Theorem~3
\cite[pp.~217--219]{Mack1993}. This section names the $\sigma$-algebra and proves that sentence.

\begin{definition}[The observed data]\label{def:obsSigma}\lean{VerifiedReserving.RandomTriangle.obsSigma, VerifiedReserving.RandomTriangle.otherObs}\leanok
\uses{def:rowSigma}
After $n$ calendar years the triangle holds one diagonal per accident year: row $i$ is observed up
to development year $n-1-i$. So
\[
  \mathrm{obsSigma} = \bigvee_{i < n} \sigma(C_{i,0},\dots,C_{i,n-1-i}),
\]
and $\mathrm{otherObs}_i$ is the same join over the accident years $i' \ne i$.
\end{definition}

\begin{lemma}[The observed data splits at one accident year]\label{lem:obsSplit}\lean{VerifiedReserving.RandomTriangle.rowSigma_le_obsSigma, VerifiedReserving.RandomTriangle.rowSigma_le_otherObs, VerifiedReserving.RandomTriangle.obsSigma_le, VerifiedReserving.RandomTriangle.otherObs_le_otherRowsAll, VerifiedReserving.RandomTriangle.otherObs_le, VerifiedReserving.RandomTriangle.obsSigma_eq_sup}\leanok
\uses{def:obsSigma}
For $i < n$, $\mathrm{obsSigma} = \sigma(C_{i,0},\dots,C_{i,n-1-i}) \vee \mathrm{otherObs}_i$, and
$\mathrm{otherObs}_i$ is contained in the $\sigma$-algebra of all the other accident years.
\end{lemma}
\begin{proof}\leanok
Both inclusions termwise: a term of the join either is row $i$'s or is one of the others'. This is
Lemma~\ref{lem:indep_rows} for the latest diagonal instead of a development year.
\end{proof}

\begin{lemma}[The chain-ladder estimators are functions of the observed data]\label{lem:clObs}\lean{VerifiedReserving.RandomTriangle.stronglyMeasurable_obsSigma_C, VerifiedReserving.RandomTriangle.stronglyMeasurable_obsSigma_Srv, VerifiedReserving.RandomTriangle.stronglyMeasurable_obsSigma_fhatRv, VerifiedReserving.RandomTriangle.stronglyMeasurable_obsSigma_ChatRv}\leanok
\uses{def:obsSigma, def:rv, def:ChatRv}
$S_k$, $\hat f_k$ and $\hat C_{i,d+m}$ are $\mathrm{obsSigma}$-measurable, with no hypothesis: an
accident year contributes to $\hat f_k$ only when $C_{i,k+1}$ lies on or above the latest diagonal.
This is the hypothesis $\hat C_{i,d+m}$ is $\mathcal{D}$-measurable of Theorem~\ref{thm:msep}.
\end{lemma}
\begin{proof}\leanok
$i \in \mathrm{contributors}(n,k)$ gives $i + k + 2 \le n$, hence $k+1 \le n-1-i$; then induction on
$m$ for the product.
\end{proof}

\begin{theorem}[No further information about one accident year]\label{thm:noFurtherInfo}\lean{VerifiedReserving.indep_rowSigmaAll_otherObs, VerifiedReserving.condExp_obsSigma_eq_rowSigma, VerifiedReserving.condExp_D_eq_rowSigma, VerifiedReserving.condExp_obsSigma_eq_D, VerifiedReserving.condExp_C_obsSigma_eq_D, VerifiedReserving.condExp_sq_C_obsSigma_eq_D}\leanok
\uses{lem:obsSplit, thm:condExp_sup, def:rowAssumptions}
Assume (M2) and that $\mathcal D_k$ is generated by the row histories. For $g$
integrable and measurable for accident year $i$, and $i < n$,
\[
  E[g \mid \mathrm{obsSigma}] = E[g \mid \sigma(C_{i,0},\dots,C_{i,n-1-i})]
  = E[g \mid \mathcal{D}_{n-1-i}] .
\]
Taking $g = C_{i,k}$ and $g = C_{i,k}^2$ gives exactly the two ``no extra information''
hypotheses of Theorem~\ref{thm:msep}.
\end{theorem}
\begin{proof}\leanok
Split the conditioning $\sigma$-algebra by Lemma~\ref{lem:obsSplit} (for the first equality) and by
Lemma~\ref{lem:indep_rows} (for the second), and apply Theorem~\ref{thm:condExp_sup} with $m_1$ row
$i$'s observed history, $m_1'$ the whole of row $i$ and $m_2$ the other rows, which (M2) makes
independent of $m_1'$. Both sides reduce to the same conditional expectation.
\end{proof}

\begin{theorem}[Mack's Theorem~3 from Mack's own assumptions]\label{thm:msepRows}\lean{VerifiedReserving.condMsep_eq_of_rows}\leanok
\uses{thm:noFurtherInfo, lem:clObs, thm:msep, thm:passage}
Assume (M1row), (M3row), (M2), that $\mathcal{D}_k$ is generated by the rows' histories, and square
integrability of the claims and of the chain-ladder prediction. Then, conditionally on the observed
data,
\[
  E\big[(\hat C_{i,d+m} - C_{i,d+m})^2 \mid \mathrm{obsSigma}\big]
  = V_m + \big(\hat C_{i,d+m} - C_{i,d}\textstyle\prod_{l<m} f_{d+l}\big)^2 ,
  \qquad d = n-1-i .
\]
Nothing is left over: the two hypotheses of Theorem~\ref{thm:msep} are
Theorem~\ref{thm:noFurtherInfo}, the measurability of the predictor is Lemma~\ref{lem:clObs}, and
(M1), (M3) in the $\mathcal{D}_k$ form are Theorem~\ref{thm:passage}.
\end{theorem}
\begin{proof}\leanok Substitute those four into Theorem~\ref{thm:msep}. \end{proof}

\begin{theorem}[The cross-term condition from independence]\label{thm:crossFree}\lean{VerifiedReserving.condExp_cross_obsSigma_eq_zero, VerifiedReserving.condCrossFree_of_rows}\leanok
\uses{thm:noFurtherInfo, def:crossFree, thm:mack2prime}
Assume (M2) and square integrability of the claims. For $i \ne j$ with $i < n$,
\[
  E\big[(C_i - E[C_i \mid \mathrm{obsSigma}])(C_j - E[C_j \mid \mathrm{obsSigma}])
  \;\big|\; \mathrm{obsSigma}\big] = 0 ,
\]
so Definition~\ref{def:crossFree} holds for the true ultimates over all $n$ accident years.
\end{theorem}
\begin{proof}\leanok
The argument of Theorem~\ref{thm:mack2prime}. Condition on the observed data together with the whole
of accident year $j$: row $j$'s centred ultimate is measurable there and comes out of the
conditional expectation. What remains is row $i$'s centred ultimate conditioned on a
$\sigma$-algebra that adds nothing about row $i$ (Theorem~\ref{thm:condExp_sup}), and the centring
was by exactly that conditional expectation (Theorem~\ref{thm:noFurtherInfo}), so it vanishes. The
tower property brings the identity back to the observed data.
\end{proof}

\begin{theorem}[Mack's Corollary from Mack's own assumptions]\label{thm:msepTotalRows}\lean{VerifiedReserving.condMsepTotal_eq_of_rows, VerifiedReserving.sum_procVar_le_condMsepTotal_of_rows}\leanok
\uses{thm:crossFree, thm:noFurtherInfo, lem:clObs, thm:msepTotalExact, thm:passage}
Under the same hypotheses as Theorem~\ref{thm:msepRows}, taken over all $n$ accident years,
\[
  E\Big[\Big(\sum_i \hat C_i - \sum_i C_i\Big)^2 \;\Big|\; \mathrm{obsSigma}\Big]
  = \sum_i \mathrm{procVar}_i + \Big(\sum_i \big(\hat C_i - M_i\big)\Big)^2
  \;\ge\; \sum_i \mathrm{procVar}_i .
\]
The chain from Mack's row-conditioned assumptions plus independence across accident years to the
exact Theorem~3 and its total-reserve form is closed.
\end{theorem}
\begin{proof}\leanok
Theorem~\ref{thm:msepTotalExact} with $\mathcal{D} = \mathrm{obsSigma}$, its cross-term hypothesis
supplied by Theorem~\ref{thm:crossFree} and its remaining hypotheses by
Theorem~\ref{thm:noFurtherInfo} and Lemma~\ref{lem:clObs}.
\end{proof}
